% Why the two broad phases differ. Left: sorting on one axis prunes on one axis,
% which a draped surface defeats. Right: the cell list bins instead, and settles
% the duplicates that binning creates with a rule that carries no per-pair state.
%
% Labels are placed beside the marks they name rather than above or below them,
% so nothing crosses a rule or a box edge, and the legend keeps the three marks
% of the right-hand panel off the drawing.
\begin{tikzpicture}[
  font=\footnotesize,
  >={Stealth[length=2mm]},
  bar/.style={draw=sccdInk, fill=sccdInk!7, line width=0.35pt},
  tested/.style={draw=sccdInk, fill=sccdAmber!18, line width=0.35pt},
  hd/.style={font=\footnotesize\itshape, text=sccdInk, inner sep=2pt},
  tag/.style={font=\scriptsize\itshape, inner sep=2pt},
  note/.style={font=\scriptsize, text=sccdInk, align=left, text width=63mm,
               anchor=north west, inner sep=0pt},
]

% ======================= left: sweep and prune ==========================
\begin{scope}
  \node[hd, anchor=south west] at (0,4.32) {sweep and prune along axis $a$};

  % Boxes sorted by their minimum on a, one per row. Many intervals overlap the
  % query's, which is what a draped surface does to any single axis.
  \foreach \i/\s/\w/\st in {
      0/0.00/1.7/bar, 1/0.35/1.9/tested, 2/0.55/2.1/tested,
      3/0.80/1.6/tested, 4/1.05/2.3/tested, 5/1.30/1.8/tested,
      6/2.95/1.7/bar,  7/3.30/1.5/bar} {
    \pgfmathsetmacro{\y}{3.70 - 0.30*\i}
    \draw[\st] (\s,\y) rectangle ++(\w,0.20);
  }

  % The query box, well clear of the lowest bar above it.
  \draw[draw=sccdAmber, fill=sccdAmber!45, line width=0.8pt]
    (0.60,1.08) rectangle ++(2.00,0.22);
  \node[tag, text=sccdAmber, anchor=west] at (2.78,1.19) {query box};

  % Its extent on a, and the window that extent opens. The span label sits
  % beside the arrow, so neither the arrow nor the axis runs through text.
  \draw[densely dashed, draw=sccdAmber] (0.60,1.06) -- (0.60,0.84);
  \draw[densely dashed, draw=sccdAmber] (2.60,1.06) -- (2.60,0.84);
  \draw[<->, draw=sccdAmber, line width=0.7pt] (0.60,0.76) -- (2.60,0.76);
  \node[tag, text=sccdAmber, anchor=west] at (2.78,0.76) {window on $a$};

  \draw[->, draw=sccdInk] (-0.05,0.34) -- (5.05,0.34)
    node[anchor=west, inner sep=3pt] {$a$};

  \node[note] at (-0.05,0.06)
    {Every shaded box overlaps the query on $a$ and is therefore tested in full.
     On a draped surface almost all of them then fail on one of the other two
     axes: about $2100$ candidates are walked for each pair actually found, and
     no choice of axis rescues it. The two good axes give $2882$ and $2838$
     candidates per pair, the third $126\,063$.};
\end{scope}

% ========================= right: 2D cell list ==========================
\begin{scope}[xshift=74mm]
  \node[hd, anchor=south west] at (0,4.32) {two-dimensional cell list};

  \draw[step=0.62, draw=sccdGrid, line width=0.4pt] (0,1.70) grid (4.96,4.18);
  \draw[draw=sccdGrid, line width=0.4pt] (0,1.70) rectangle (4.96,4.18);

  % Teal and plum are the two primitives throughout the paper: the same two
  % whose difference function Equations (1) and (2) root.
  \draw[draw=sccdTeal, fill=sccdTeal!14, line width=0.7pt]
    (0.62,1.95) rectangle (2.80,3.55);
  \node[tag, text=sccdTeal, anchor=south west] at (0.66,1.99) {first};
  \draw[draw=sccdPlum, fill=sccdPlum!14, line width=0.7pt]
    (2.10,3.10) rectangle (4.30,3.95);
  \node[tag, text=sccdPlum, anchor=north east] at (4.26,3.91) {second};

  % The overlap, then the owning cell over both fills. The cell is outlined in
  % ink and dashed so it stays legible where it crosses the amber overlap.
  \fill[sccdAmber!55] (2.10,3.10) rectangle (2.80,3.55);
  \draw[draw=sccdInk, line width=1.1pt, densely dashed]
    (1.86,2.94) rectangle ++(0.62,0.62);
  \fill[sccdInk] (2.10,3.10) circle (1.9pt);

  % --- legend, one item per line so nothing collides at this size ---
  \fill[sccdAmber!55] (0.05,1.34) rectangle ++(0.28,0.20);
  \node[font=\scriptsize, anchor=west, inner sep=3pt] at (0.33,1.44)
    {the overlap of the two primitives' boxes};
  \draw[draw=sccdInk, line width=1.1pt, densely dashed]
    (0.05,1.02) rectangle ++(0.28,0.20);
  \node[font=\scriptsize, anchor=west, inner sep=3pt] at (0.33,1.12)
    {the cell that owns the pair};
  \fill[sccdInk] (0.19,0.80) circle (1.9pt);
  \node[font=\scriptsize, anchor=west, inner sep=3pt] at (0.33,0.80)
    {the overlap's minimum corner};

  \node[note] at (-0.05,0.56)
    {Each box is binned into every cell it touches, so the same pair can be met
     in several of them. It is emitted only by the cell holding the minimum
     corner of the overlap. Since that corner lies inside both boxes, both are
     binned there and the cell is unique: two clamps per surviving pair, and no
     per-pair state to allocate.};
\end{scope}
\end{tikzpicture}
